\documentclass[a4paper,11pt]{ltjsarticle}
\usepackage[margin=24mm,headheight=16pt]{geometry}
\usepackage{amsmath,amssymb,amsthm,mathtools,mathrsfs}
\usepackage{booktabs,longtable,array,enumitem,needspace}
\usepackage[unicode,colorlinks=true,linkcolor=blue!45!black,urlcolor=blue!45!black,citecolor=blue!45!black]{hyperref}
\usepackage{xcolor,fancyhdr}
\pagestyle{fancy}\fancyhf{}\fancyhead[L]{\small Poisson各成分の支配的chart}\fancyhead[R]{\small 2026-11-09 / v11}\fancyfoot[C]{\thepage}
\newcommand{\CC}{\mathbb C}\newcommand{\PP}{\mathbb P}\newcommand{\OO}{\mathcal O}
\newcommand{\Perf}{\operatorname{Perf}}\newcommand{\Ku}{\operatorname{Ku}}
\newtheorem{theorem}{定理}[section]\newtheorem{proposition}[theorem]{命題}
\newtheorem{lemma}[theorem]{補題}\newtheorem{corollary}[theorem]{系}
\theoremstyle{definition}\newtheorem{definition}[theorem]{定義}\newtheorem{example}[theorem]{例}
\theoremstyle{remark}\newtheorem{remark}[theorem]{注意}
\renewcommand{\proofname}{証明}\allowdisplaybreaks
\setlength{\parskip}{0.15em}
\title{Poisson既約成分の支配的chartと可換復元\\\large 中心三次元の線形代数化とsaturationの有限証明書}
\author{}\date{2026年11月9日・v11\\\small 論文Cの現行節：独立読解用草稿}
\begin{document}\maketitle
\begin{abstract}
任意次元の滑らかなcubicのPoissonスキームの各既約成分について、
特異な零点上で普遍量子化を構成することを前提とせず、
滑らかな支配的なスカラーchartでgenericな可換復元を研究する方法を与える。
必要なら有限étale被覆上でブロックを固定する。
有限行列の形式核ベクトルは、任意の有限jetを保って代数化できることを証明し、中心三次元へ適用する。
残余dg圏の相対saturationと既存原稿のétale局所実現原理を仮定すれば、
一つの可換ファイバーから各成分の一般点へ可換模型を移せる。
相対saturationそのものと一般の可換復元の存在は未証明である。
\end{abstract}
\tableofcontents
\section{一般次元におけるPoisson既約成分ごとの可換復元}
\label{sec:all-poisson-components}

本節の一般問題は、固定した三次元族ではなく、任意の$n\ge3$と
滑らかなcubic $X^n$について、Poissonスキームの\emph{各既約成分}の
一般点の量子化からrotation付き可換模型を得ることである。
任意のAS-regular中心三次商を無条件に対象とする問題とは区別する。

\begin{proposition}[Poisson錐の各既約成分は可換点を含む]
\label{prop:all-components-cone}
$X$を滑らかな複素射影多様体とし、
\[
\operatorname{Pois}(X)=
\left\{\pi\in H^0(X,\bigwedge^2T_X)\mid[\pi,\pi]=0\right\}
\]
をアフィンPoissonスキームとする。
その被約化の任意の既約成分$Z$は零Poisson構造$0$を含む。
$X$が滑らかなcubicであれば、同じ主張はcubicも変化させる
相対Poissonスキームの各既約成分について成り立つ。
\end{proposition}
\begin{proof}
Jacobi方程式は$\pi$について二次斉次である。
従って$\mathbb G_m$は各Poissonスキームに$\pi\mapsto t\pi$で作用する。
連結な$\mathbb G_m$は有限個の既約成分の集合を自明に作用するので、
$Z$はこの作用で不変である。
$z=(X,\pi)\in Z$について
$\mathbb A^1\to\operatorname{Pois}$、
$t\mapsto(X,t\pi)$は零まで定義され、
$t\ne0$で像が$Z$にあるから$(X,0)\in Z$である。
固定$X$の場合も同じである。
\end{proof}

\begin{theorem}[各Poisson成分の一般点へ進むための条件付き原理]
\label{thm:all-components-conditional}
$n\ge3$とし、$Z$を滑らかなcubic $X^n$上の
アフィンPoissonスキームの被約な既約成分とする。
$0\in U\subset Z$をZariski開近傍とし、
$U$上に中心三次商の代数的な次数別flat量子化族が存在すると仮定する。
さらに、その残余dg圏と標準rotation
$(\mathcal K_U,\rho_U)$が$U$-linearなsmooth proper
代数的pair族をなし、零ファイバーが
$(\Ku(X),\rho_X)$とenhanced同値であると仮定する。
本論文のétale局所復元定理\ref{thm:etale}がこの族に適用できるなら、
$0$を含む非空開部分$U^\circ\subset U$が存在し、
各複素点におけるpairは滑らかな可換cubicのpairと同値である。
特に$Z$の幾何学的一般点は、有限分離拡大の後、
rotation付き可換cubic模型を持つ。
\end{theorem}
\begin{proof}
命題\ref{prop:all-components-cone}により$0\in Z$である。
仮定と定理\ref{thm:etale}から$0$を通るpointed étale近傍で
相対cubicとrotationを保つ同値が存在する。
その像は非空Zariski開であり、既約な$Z$で稠密である。
generic point上ではétale近傍の支配成分の剰余体が
$\mathbb C(Z)$の有限分離拡大となるので結論を得る。
\end{proof}
\begin{remark}[何が未証明か]
\label{rem:all-components-gap}
上の定理は\emph{条件付き}である。
個々のPoisson構造の一パラメータ量子化の存在から、
その成分$Z$全体の零近傍における有限型の一つのflat量子化族は
自動的には従わない。
特にブロック数の変化、中心元のchart間の接着、
相対Beilinson条件、有限対角証明書、相対Orlov比較と底変更を
一様に確立する必要がある。
また、射影化されたPoissonパラメータ空間そのものには零点は無い。
ここで用いるのはその\emph{アフィン錐}である。
\end{remark}

\begin{proposition}[次元の合同類とrotationのSerre関手からの復元]
\label{prop:all-components-rotation-serre}
$n\ge3$、$\mathcal K$をSerre関手$S_{\mathcal K}$を持つ
$k$-linear smooth proper dg圏、$\rho$を自己同値とし、
\[
 \rho^3\simeq[2],\qquad
 S_{\mathcal K}\simeq\rho^{1-n}[n]
 \tag{SR}\label{eq:all-components-serre}
\]
を仮定する。
$n\not\equiv1\pmod3$なら、$3a+(1-n)b=1$を満たす
整数$a,b$に対して
\[
 \boxed{\quad \rho\simeq S_{\mathcal K}^{\,b}[\,2a-nb\,]\quad}
 \tag{R}\label{eq:all-components-rotation}
\]
である。従ってこの範囲では、上の二恒等式を満たす
rotationの同型類はSerre関手とshiftから決まる。
$n\equiv1\pmod3$なら二恒等式からは
$S_{\mathcal K}\simeq[(n+2)/3]$しか導かれず、
$\rho$の同型類は一般に決定されない。
\end{proposition}
\begin{proof}
前半は$\gcd(3,n-1)=1$によるBézout等式から
$\rho=(\rho^3)^a(\rho^{1-n})^b
\simeq[2a]S_{\mathcal K}^{\,b}[-nb]$
を得る。
$n\equiv1\pmod3$なら$n-1=3m$なので
$S_{\mathcal K}\simeq(\rho^3)^{-m}[n]
\simeq[n-2m]=[(n+2)/3]$となる。
この式だけから$\rho$を復元できないことは、
本論文のcubic fourfoldに対するZhang twistのrotation障害の例が示す。
\end{proof}
\begin{remark}[全非可換成分との区別]
\label{rem:all-components-nc-scope}
任意のAS-regular中心三次商\emph{全ての点}の
rotation付き可換復元は本論文のZhang twist反例で偽である。
ただし個別の反例点だけで、非可換中心三次対の任意の
既約成分の\emph{一般点}における可換復元を否定したことにはならない。
Poisson零点を含む各成分の問題と、
可換点に接続しない非可換モジュライ成分の問題は別に扱う。
$n=1,2$ではKuznetsov残余圏とPoisson変形問題の範囲が変わるため、
本節の$n\ge3$の条件付き原理をそのまま拡張しない。
\end{remark}

\begin{thebibliography}{99}
\bibitem{Stacks} The Stacks Project Authors, \emph{The Stacks Project}.
Tags 00MC (completion and faithful flatness), 054K (Chevalley),
0BDK (perfect-complex vanishing opens), 055U (étale-local sections of smooth morphisms).
\url{https://stacks.math.columbia.edu/tag/00MC};
\url{https://stacks.math.columbia.edu/tag/0BDK};
\url{https://stacks.math.columbia.edu/tag/055U}.
\bibitem{TV} B. Toën and M. Vaquié,
\emph{Moduli of objects in dg-categories},
Ann. Sci. École Norm. Sup. (4) \textbf{40} (2007), 387--444;
arXiv:math/0503269v5, Theorem 3.6.
\url{https://arxiv.org/abs/math/0503269}.
\end{thebibliography}
\end{document}
